\chapter{THỰC NGHIỆM VÀ ĐÁNH GIÁ}
\label{ch:thuc-nghiem}

\section{Giới thiệu}
\label{sec:gioi-thieu-c4}
% Nêu mục đích chương: hiện thực hóa phương pháp thành công cụ và đánh giá theo bốn câu hỏi nghiên cứu.

\section{Công cụ hỗ trợ}
\label{sec:cong-cu}
% Kiến trúc công cụ ddd-codegen; ánh xạ năm bước phương pháp sang thành phần công cụ; quy mô cài đặt.

\section{Ca nghiên cứu CourseMan}
\label{sec:ca-nghien-cuu}
% Mô tả hai ngữ cảnh giới hạn, chính sách kiểm soát truy cập, thống kê đặc tả đầu vào.

\section{Thiết kế đánh giá}
\label{sec:thiet-ke-danh-gia}
% Bốn câu hỏi nghiên cứu, thước đo, nguồn dữ liệu; quy trình bốn bước; bộ kịch bản kiểm thử.

\section{Kết quả thực nghiệm}
\label{sec:ket-qua}
% Kết quả theo RQ1--RQ4.

\section{Các nguy cơ ảnh hưởng tính hợp lệ}
\label{sec:threats}
% Internal, external, construct, conclusion validity và biện pháp giảm thiểu.

\section{Bàn luận}
\label{sec:ban-luan}
% Đối sánh với công trình liên quan; bài học rút ra; giới hạn của công cụ.

\section{Tổng kết chương}
\label{sec:tong-ket-c4}
